\section*{第 \ref{ch:b3:innate-immunity} 章——先天免疫与炎症}

\begin{solution}{exo:b3:innate-immunity:1}
\omterm{def:b3:innate-immunity:innate}{屏障}包括：皮肤角蛋白（机械\omterm{def:b3:innate-immunity:innate}{屏障}）、气道黏液与\omterm{def:b3:cytoskeleton-motility:cilia}{纤毛}（捕获
并清除）、胃酸（杀灭吞下的微生物）、分泌液中的溶菌酶与
防御素（裂解并穿透细胞膜）、\omterm{def:b3:microbiomes:symbiosis}{共生}\omterm{def:b3:microbiomes:symbiosis}{微生物组}
（占据生态位）。细胞包括：\omterm{def:b3:innate-immunity:innate}{中性粒细胞}（通过\omterm{def:b3:innate-immunity:killing}{吞噬作用}杀灭细菌）、
\omterm{def:b3:innate-immunity:innate}{巨噬细胞}（吞噬、清理并发出警报）、\omterm{def:b3:innate-immunity:innate}{树突状细胞}（将
抗原运送至淋巴结并启动适应性应答）、肥大细胞
（释放组胺、引发\omterm{def:b3:innate-immunity:inflammation}{炎症}）、\omterm{def:b3:innate-immunity:innate}{自然杀伤细胞}（杀灭被感染的
或已转化的细胞）。
\end{solution}

\begin{solution}{exo:b3:innate-immunity:2}
为一整类微生物所共有，对它们必不可少（因此无法通过丢失来逃避识别），
且宿主体内不存在。\omterm{def:b3:bacteriology:envelope}{脂多糖}——TLR4；
鞭毛蛋白——TLR5；双链 RNA——TLR3（胞质中还有 RIG-I）；
未甲基化的 \omterm{def:b3:chromatin-epigenetics:methylation}{CpG} DNA——TLR9；\omterm{def:b3:bacteriology:envelope}{肽聚糖}片段——NOD
蛋白；$\beta$-葡聚糖——dectin-1。
\end{solution}

\begin{solution}{exo:b3:innate-immunity:3}
发红源于组胺、一氧化氮和
前列腺素引起的小动脉扩张；\omterm{prop:b3:innate-immunity:systemic}{发热}源于同样增加的血流。肿胀：
内皮收缩后小静脉\omterm{def:b3:innate-immunity:inflammation}{渗出}血浆（由组胺、C3a、C5a 引起）。
疼痛：缓激肽和前列腺素 E$_{2}$ 使伤害感受器敏化，
外加肿胀所致的压迫。
\end{solution}

\begin{solution}{exo:b3:innate-immunity:4}
经典途径：C1q 结合表面上的抗体（或 C 反应蛋白）。
凝集素途径：甘露糖结合凝集素结合微生物糖类。旁路途径：
C3 的自发缓慢活化及 C3b 在任何无保护表面上的沉积。
结果：C3b 作为调理素促进吞噬；C3a 和 C5a 引发\omterm{def:b3:innate-immunity:inflammation}{炎症}并
募集细胞；C5b 启动形成裂解细胞的\omterm{def:b3:innate-immunity:complement}{膜攻击复合物}。
\end{solution}

\begin{solution}{exo:b3:innate-immunity:5}
$B = [k_{0}/(a-d)](e^{(a-d)t} - 1)$，其中 $a - d = 0.1$；在 \qty{60}{s} 时，
$20(e^{6} - 1) \approx 8000$；在 \qty{90}{s} 时，$20(e^{9} - 1) \approx
1.6\times 10^{5}$。宿主细胞：$k_{0}/(d - a) = 2/0.1 = 20$。
\end{solution}

\begin{solution}{exo:b3:innate-immunity:6}
沿\omterm{def:b3:innate-immunity:inflammation}{选择素}滚动（凝集素与糖类结合）；\omterm{def:b3:innate-immunity:inflammation}{趋化因子}
信号激活\omterm{def:b3:innate-immunity:inflammation}{整合素}；\omterm{def:b3:innate-immunity:inflammation}{整合素}与内皮黏附分子牢固黏附
（ICAM）；穿越内皮细胞之间的间隙迁移
；沿梯度趋化。若缺乏
$\beta_{2}$ \omterm{def:b3:innate-immunity:inflammation}{整合素}，\omterm{def:b3:innate-immunity:innate}{中性粒细胞}只能滚动而不能停止或穿越：
它们在血液中堆积（白细胞计数极高），感染反复发生而不形成
脓液，伤口愈合不良，脐带脱落延迟。
\end{solution}

\begin{solution}{exo:b3:innate-immunity:7}
慢性肉芽肿病：NADPH 氧化酶有缺陷，因此吞噬体中不产生
超氧阴离子、过氧化氢或次氯酸盐。
由能分解自身过氧化氢的生物所引起的感染
（过氧化氢酶阳性）：\emph{Staphylococcus aureus}、\emph{Aspergillus}、
\emph{Burkholderia}、\emph{Serratia}、\emph{Nocardia}。肉芽肿形成的原因是
\omterm{def:b3:innate-immunity:innate}{巨噬细胞}吞噬了无法杀灭的微生物，并将其
包裹隔离，在慢性\omterm{def:b3:dna-repair:lesions}{损伤}中不断有细胞募集。
\end{solution}

\begin{solution}{exo:b3:innate-immunity:8}
抑制性受体检测 MHC~I；激活性受体检测应激配体，
若激活超过抑制，细胞即死亡。(a) 丢失
MHC~I 解除了抑制：\omterm{def:b3:innate-immunity:innate}{自然杀伤细胞}杀伤该细胞——这正是
逃避 T 细胞所付出的代价。(b) 若保留 MHC~I 但表达应激配体，
激活仍可占优势，细胞同样被杀伤；因此\omterm{def:b3:cancer-biology:hallmarks}{肿瘤}
常会脱落或抑制应激配体。
\end{solution}

\begin{solution}{exo:b3:innate-immunity:9}
纯蛋白质不携带病原体模式，因此摄取它的\omterm{def:b3:innate-immunity:innate}{树突状细胞}
不会成熟，无共刺激地呈递它，而识别它的
T 细胞要么被关闭，要么不予理会。\omterm{prop:b3:innate-immunity:bridge}{佐剂}
——明矾、脂质、Toll 样受体配体——可触发\omterm{def:b3:innate-immunity:prr}{模式
识别受体}，使\omterm{def:b3:innate-immunity:innate}{树突状细胞}成熟，并提供第二信号。
这正是贾内韦的观点：适应性系统只有在
先天系统判定抗原随微生物而来时，才对其作出应答。
\end{solution}

\begin{solution}{exo:b3:innate-immunity:10}
无\omterm{prop:b3:innate-immunity:systemic}{发热}时，经 $48$ 次倍增：$10^{4}\times 2^{48} \approx 3\times
10^{18}$（一个荒谬的结果，说明决定结局的是杀灭而非
生长）。有\omterm{prop:b3:innate-immunity:systemic}{发热}时，经 $32$ 次倍增得 $4\times 10^{13}$，
\omterm{def:b3:innate-immunity:innate}{中性粒细胞}需面对的细菌数因此减少约 $2^{16} \approx 65\,000$ 倍。
代价：$3\times 0.12\times 8 = \qty{2.9}{MJ}$ 每天，相当于一顿大餐。
当病原体生长被减缓且宿主有储备时值得；
若病原体在 \qty{40}{\celsius} 下同样生长，或宿主
正在挨饿，或\omterm{prop:b3:innate-immunity:systemic}{发热}本身接近危险温度时，则不值得。
\end{solution}

\begin{solution}{exo:b3:innate-immunity:11}
结合因子 H 招募宿主自身的调节因子：$d$ 升高到
超过 $a$，循环在微生物表面与在宿主细胞上一样衰减。裂解
C3b 可清除已沉积的转化酶，同样使 $d$ 升高。荚膜
呈现的表面上转化酶难以形成，从而降低 $a$，并
遮蔽吞噬细胞受体所结合的 C3b。最省的办法是：
一种结合因子 H 的表面蛋白，宿主免费提供因子 H，
由其发挥作用。
\end{solution}

\begin{solution}{exo:b3:innate-immunity:12}
在刺伤处，扩张、渗漏与募集只涉及
几毫升范围内的血管，并在需要之处提供防御；同样的
介质若释放到整个循环系统中，会使每一条血管扩张、
每一根毛细血管渗漏，因此血压崩溃，凝血
随处被触发——生理机制相同，但规模
足以致命。TNF 驱动类风湿关节炎的滑膜\omterm{def:b3:innate-immunity:inflammation}{炎症}，
因此阻断它可缓解该病；但 TNF 也维持
肉芽肿中\omterm{def:b3:innate-immunity:innate}{巨噬细胞}的组织结构，这些肉芽肿包裹着休眠的结核杆菌，
阻断它会使潜伏性结核得以逃逸。
\end{solution}

\begin{solution}{pb:b3:innate-immunity:1}
\textbf{1.} 一小时相当于 $1.5$ 次倍增：$10^{4}\times 2^{1.5} \approx
2.8\times 10^{4}$.
\textbf{2.} $2\times 10^{4}$ 个\omterm{def:b3:innate-immunity:innate}{中性粒细胞}最多可杀灭 $4\times 10^{5}$ 个；
细菌仅从 $2.8\times 10^{4}$ 增长至 $8\times 10^{4}$：
杀灭量超过增殖量五倍。
\textbf{3.} 第 2 小时：杀灭前为 $8\times 10^{4}$，杀灭后为零。
种群在 \qty{2}{h} 附近达到约 $8\times 10^{4}$ 的峰值，
此后为零；至 \qty{6}{h} 时已无剩余。
\textbf{4.} \omterm{def:b3:innate-immunity:innate}{中性粒细胞}于 \qty{3}{h} 到达，此时细菌数为 $10^{4}\times 2^{4.5} = 2.3\times
10^{5}$。第 4 小时：$6.4\times 10^{5} - 4\times 10^{5} = 2.4\times
10^{5}$；第 5 小时：$6.8\times 10^{5} - 4\times 10^{5} = 2.8\times 10^{5}$；
第 6 小时：$7.9\times 10^{5} - 4\times 10^{5} = 3.9\times 10^{5}$，
仍在上升——\omterm{def:b3:innate-immunity:innate}{中性粒细胞}已无法跟上，于是形成脓肿。两
小时的延迟，便会使可治愈的感染转为慢性感染。
\textbf{5.} $8\times 10^{4}/20 = 4000$ 个\omterm{def:b3:innate-immunity:innate}{中性粒细胞}因杀菌而死亡，另加
$2\times 10^{4}$ 个到达后一天内无论如何都会死亡的细胞。脓液由
死亡的\omterm{def:b3:innate-immunity:innate}{中性粒细胞}、死亡的细菌及液化组织构成；\omterm{def:b3:innate-immunity:innate}{巨噬细胞}
通过吞噬尸体将其清除，否则必须引流。
\textbf{6.} \omterm{def:b3:innate-immunity:innate}{中性粒细胞}每小时新到达 $2000$ 个，每小时杀灭 $4\times 10^{4}$ 个：第 2 小时，
$8\times 10^{4} - 4\times 10^{4} = 4\times 10^{4}$；第 3 小时，$1.1\times
10^{5} - 4\times 10^{4} = 7\times 10^{4}$；第 4 小时，$2\times 10^{5} -
4\times 10^{4} = 1.6\times 10^{5}$——无限增长。
患者无法遏制如此微小的接种量，因此\omterm{def:b3:bacteriology:antibiotics}{抗生素}须
自最初迹象出现起即承担杀菌任务。
\textbf{7.} $B = 12.5(e^{0.08t} - 1)$：当 \qty{30}{s} 时为 $125$；\qty{60}{s} 时
为 $1500$；\qty{120}{s} 时为 $1.8\times 10^{5}$。
\textbf{8.} $e^{0.08t} = 1.6\times 10^{5}$：$t = 12/0.08 = \qty{150}{s}$。
\textbf{9.} $k_{0}/(d - a) = 10$.
\textbf{10.} $B(120) = 33(e^{3.6} - 1) \approx 1200$，少 $150$ 倍。
荚膜为细菌争取了数分钟——既是增殖的时间，也是保护，直至
抗体到达，以便在荚膜本身上启动
经典途径。
\textbf{11.} $10^{4}\times 2\times 10^{6} = 2\times 10^{10}$ 个 C3b，
即血浆中 C3 的百万分之一。
\textbf{12.} 若缺乏 C3，三条途径都将停在其共同步骤处：
无\omterm{def:b3:innate-immunity:complement}{调理作用}、无\omterm{def:b3:innate-immunity:complement}{过敏毒素}、无裂解——严重的反复
荚膜细菌感染。若缺乏 C9，则仅
失去末端孔道；\omterm{def:b3:innate-immunity:complement}{调理作用}与\omterm{def:b3:innate-immunity:inflammation}{炎症}均正常，表型为
反复的 \emph{Neisseria} 感染——这是\omterm{def:b3:innate-immunity:killing}{吞噬作用}控制不力、而裂解作用控制良好的
唯一一个属。
\textbf{13.} $2\times 0.24\times 8 = \qty{3.8}{MJ}$，约合 \qty{100}{g}
脂肪。
\textbf{14.} 一小时内仅一次倍增：$2\times 10^{4}$，而非
$2.8\times 10^{4}$，相差 $1.4$ 倍。
\textbf{15.} 第 2 小时：杀灭前为 $4\times 10^{4}$，杀灭后为零——至
\qty{2}{h} 时已清除，与前相同；\omterm{prop:b3:innate-immunity:systemic}{发热}的作用在
\omterm{def:b3:innate-immunity:innate}{中性粒细胞}供应不足时尤为重要，如第 4 题和第 6 题所示。
\textbf{16.} $199\,\text{mg/L}\times 3\,\text{L} = \qty{0.6}{g}$；$0.6/
\num{115000} = 5.2\times 10^{-6}$ 摩尔 $= 3.1\times 10^{18}$ 个分子；
在 \qty{172800}{s} 内，即每秒 $1.8\times 10^{13}$ 个。
\textbf{17.} 它们使调定点回降至接近 \qty{37}{\celsius}，
患者感觉好转，细菌倍增略快；在此
模型中\omterm{def:b3:innate-immunity:innate}{中性粒细胞}仍能取胜，而认为退热药会加重普通感染的证据
很薄弱。
\textbf{18.} 高于约 \qty{41}{\celsius} 时，蛋白质开始变性，
酶失效，神经元出现误放电（惊厥）；\qty{42}{\celsius} 即为
致死温度。调定点低于此值，靠的是降温因子——白细胞介素-10，
加压素、糖皮质激素——以及下丘脑中前列腺素作用的
上限。
\textbf{19.} $10^{9}\times 3\times 10^{6} = 3\times 10^{15}$ 个分子
$= 5\times 10^{-9}$ 摩尔，置于 \qty{5}{L} 中：\qty{1e-9}{mol/L}。
\textbf{20.} 等于饱和浓度：体内每个\omterm{def:b3:innate-immunity:innate}{巨噬细胞}
均同时被最大程度地激活。
\textbf{21.} $10^{11}\times 1000\times 3600 = 3.6\times 10^{17}$ 个分子
$= 6\times 10^{-7}$ 摩尔；在 \qty{15}{L} 中为 \qty{4e-8}{mol/L}——约为
活性浓度的四千倍。
\textbf{22.} 每条小动脉扩张，每条小静脉渗漏：血液
血容量淤积于扩张的血管中并渗入组织，因此
血压下降；肾脏灌注不足，停止过滤；并且
在数千根毛细血管中引发的凝血将其堵塞，并消耗
凝血因子。
\textbf{23.} 被杀死的\omterm{def:b3:bacteriology:envelope}{革兰氏阴性}菌会一次性释放其全部\omterm{def:b3:bacteriology:envelope}{脂多糖}，
形成对 TLR4 的大剂量刺激，因此\omterm{def:b3:innate-immunity:inflammation}{细胞因子}激增可持续数小时而加重。
抗 TNF 疗法失效，是因为休克出现时级联反应已越过 TNF
——白细胞介素-1、白细胞介素-6 与凝血已在进行
——且控制细菌仍需要 TNF：问题在于时机，而不在
靶点本身。
\textbf{24.} 注射\omterm{def:b3:bacteriology:envelope}{脂多糖}对 TLR4 缺失
小鼠毫无作用，该小鼠可耐受对正常小鼠致死的剂量；活的\omterm{def:b3:bacteriology:envelope}{革兰氏阴性}
感染却会使其死亡，因为它无法及早察觉细菌，
也就无法发起控制细菌的\omterm{def:b3:innate-immunity:inflammation}{炎症}。悖论由此消解：引发休克的
识别分子，恰恰是提供
防御的识别分子。
\textbf{25.} 至 \qty{6}{h} 时已无细菌残留（均在 \qty{2}{h} 前被杀灭）；
\qty{120}{s} 时每个细菌约带 $1.8\times 10^{5}$ 个 C3b；两天的
\omterm{prop:b3:innate-immunity:systemic}{发热}消耗 \qty{3.8}{MJ}，约合 \qty{100}{g} 脂肪。
\end{solution}
